\chapter{Multi-Agent Topologies}
\label{ch:multi-agent}

\epigraph{``There is no spoon.''}{Spoon Boy, \emph{The Matrix} (1999)}

The line is usually quoted as mysticism. It is actually a debugging technique:
the constraint you are struggling against is a category you invented. In
multi-agent systems the invented category is ``agent'' as a special kind of
thing. There is no agent. There is a host that also happens to be a server, and
once you see that, the whole topic collapses into things you already know.

\section{An agent is a host and a server at the same time}

That is the entire architectural insight.

\begin{itemize}[leftmargin=1.6em, itemsep=2pt]
\item On its \textbf{consumer} side, it is a host: it connects to servers,
merges catalogues, runs a loop.
\item On its \textbf{provider} side, it is a server: it exposes tools that
happen to be implemented by running a loop.
\end{itemize}

Nothing in the protocol distinguishes ``a tool that queries a database'' from ``a
tool that runs a four-step agent loop and returns a synthesis''. The caller sees
a tool call. This is a good property, and it is why MCP does not need a
multi-agent extension.

\begin{py}
risk_agent = Server("meridian-risk-agent", "1.0.0", instructions=(
    "Delegated credit analysis. Give it an account and a question; it will "
    "consult the risk, fraud, and market-data servers and return a synthesis."
))

@risk_agent.tool(
    "analyse_account",
    "Full credit analysis of one account: risk score, fraud signals, and "
    "indicative pricing, synthesised into a recommendation.",
    {
        "type": "object",
        "properties": {
            "accountId": {"type": "string", "pattern": "^ACC-[0-9]{4}$"},
            "question": {"type": "string",
                         "description": "What the caller wants to know"},
        },
        "required": ["accountId"],
    },
)
def analyse_account(ctx: RequestContext):
    budget = _inherit_budget(ctx)          # see below
    result = AgentLoop(inner_host, model,
                       max_steps=budget.steps,
                       token_budget=budget.tokens,
                       cost_budget_usd=budget.dollars).run(
        f"Analyse {ctx.arguments['accountId']}. "
        f"{ctx.arguments.get('question', '')}")
    return text_result(result.answer,
                       structured={"iterations": len(result.iterations),
                                   "costUsd": result.cost_usd})
\end{py}

A tool whose implementation is a loop. The caller cannot tell, and mostly should
not need to.

\section{Three shapes}

\bookfig[1.0]{multi-agent}{The three topologies, what each one buys, and what
each one costs. The box at the bottom applies to all three.}{multi-agent}

\subsection{Supervisor tree}

One coordinator, several specialists. The shape most teams reach for, and
usually correctly.

\textbf{Isolation:} a specialist that fails, hangs, or misbehaves is contained by
the supervisor. It can retry, route around, or report.

\textbf{Failure mode:} the supervisor is a single point of failure and, more
insidiously, a single point of context. Every specialist's summary lands in the
supervisor's window, which fills.

\subsection{Peer mesh}

Agents call each other directly. No coordinator.

\textbf{Isolation:} none, structurally. Peers route around damage, which is
resilience of a different kind.

\textbf{Failure mode:} cycles. A calls B calls C calls A, and now you have an
infinite recursion distributed across three machines with a token meter attached.

\subsection{Hierarchy}

Levels, each summarising for the one above. Portfolio, sector, account.

\textbf{Isolation:} good, and the context property is the real prize: each level
sees summaries rather than raw results, so the top-level window stays small
regardless of how many accounts exist.

\textbf{Failure mode:} depth multiplies latency. Three levels means the leaf's
model turn, the middle's model turn, and the top's model turn, in series. A
1.4-second task becomes a four-second task.

\section{The four counters}

Whatever the shape, four things must travel with a delegated request, and all
four ride in \texttt{\_meta}.

\begin{py}
DEPTH_KEY   = "com.meridian/delegationDepth"
TOKENS_KEY  = "com.meridian/tokenBudgetRemaining"
DOLLARS_KEY = "com.meridian/costBudgetRemaining"
# `traceparent` is already reserved by the spec for the fourth.

MAX_DELEGATION_DEPTH = 3

def _inherit_budget(ctx: RequestContext) -> Budget:
    """Read the caller's remaining budget, or assume the smallest.

    An agent that receives no budget must not assume an unlimited one. That
    default is how a misconfigured caller turns into a four-figure invoice.
    """
    meta = ctx.raw_meta
    depth = int(meta.get(DEPTH_KEY, 0))
    if depth >= MAX_DELEGATION_DEPTH:
        raise errors.InvalidParams(
            f"Delegation depth {depth} exceeds the limit of "
            f"{MAX_DELEGATION_DEPTH}")
    return Budget(
        depth=depth + 1,
        tokens=int(meta.get(TOKENS_KEY, DEFAULT_TOKEN_BUDGET)),
        dollars=float(meta.get(DOLLARS_KEY, DEFAULT_COST_BUDGET)),
        steps=max(1, MAX_STEPS - depth * 2),
    )
\end{py}

Four decisions in that function, each worth stating:

\textbf{Depth is checked before any work.} Refusing at depth 3 costs one request.
Discovering the recursion at depth 30 costs thirty.

\textbf{The default budget is small, not infinite.} An agent that receives no
budget is talking to a caller that has not thought about budgets, which is
exactly when a conservative default matters most.

\textbf{Steps shrink with depth.} A leaf agent does not need eight iterations. If
it does, the decomposition is wrong.

\textbf{Budgets are \emph{remaining}, not allocated.} This is the one people get
wrong. Passing ``you may spend \$0.50'' to three children authorises \$1.50.
Passing what is left, and decrementing as results return, authorises \$0.50
total.

\begin{dangerbox}{Budget fan-out is multiplicative}
A supervisor with a \$1.00 budget delegating to 3 agents, each delegating to 3
more, each running 8 steps: if every level passes its full budget down, the
system is authorised to spend \$9.00 and will happily do so.

Depth limits alone do not fix this. You need the budget to be a shared,
decrementing quantity rather than a per-node allowance.
\end{dangerbox}

\section{Circuit breakers}

Budgets stop runaway spending. Circuit breakers stop runaway \emph{waiting}.

\begin{py}
class CircuitBreaker:
    """Stop calling something that is clearly broken.

    Three states. Closed: normal. Open: fail fast without calling. Half-open:
    let exactly one probe through and decide from its outcome. The half-open
    state is what stops a recovering service being flattened by every client
    reconnecting at once.
    """

    def __init__(self, *, threshold: int = 5, cooldown: float = 30.0):
        self.threshold = threshold
        self.cooldown = cooldown
        self.failures = 0
        self.opened_at: float | None = None

    def allows(self, now: float) -> bool:
        if self.opened_at is None:
            return True
        if now - self.opened_at >= self.cooldown:
            self.opened_at = None      # half-open: one probe gets through
            self.failures = self.threshold - 1
            return True
        return False

    def record(self, ok: bool, now: float) -> None:
        if ok:
            self.failures = 0
            self.opened_at = None
        else:
            self.failures += 1
            if self.failures >= self.threshold:
                self.opened_at = now
\end{py}

Failing fast is better than timing out, for a reason specific to agent systems:
a timeout consumes the caller's step budget while producing no information. A
fast failure lets the model try something else with the steps it has left.

\section{Discovery at scale}

Two servers you configure by hand. Two hundred you do not.

\begin{table}[htbp]
\centering
\small
\begin{tabularx}{\textwidth}{@{}lXX@{}}
\toprule
\thead{Mechanism} & \thead{Good for} & \thead{Watch out for} \\
\midrule
Static config (\texttt{.mcp.json})
  & Development, small fixed fleets
  & Drift between environments \\
The official MCP Registry
  & Public, discoverable servers
  & Anything public is untrusted (\chapref{threats}) \\
A private registry
  & Enterprise fleets, approval workflows
  & You now operate a registry \\
Runtime binding
  & Selecting a server per task from a catalogue
  & Discovery latency in the critical path \\
\bottomrule
\end{tabularx}
\caption{Discovery options. Most organisations end up with static config for
development and a private registry for production.}
\label{tab:discovery}
\end{table}

Whatever you use, \texttt{server/discover} is cacheable with a TTL, so honour it.
A supervisor rediscovering twelve servers per task is spending twelve round trips
on information that changes when somebody deploys.

\section{Tracing across the tree}

This is where multi-agent systems become debuggable or do not.

One trace id, propagated through every hop, so a request that fans out to three
agents and eleven servers is one trace with a shape you can read.

\begin{py}
body["_meta"] = build_request_meta(
    self.capabilities, self.info,
    protocol_version=self.protocol_version,
    progress_token=progress_token,
    traceparent=traceparent,
)
\end{py}

W3C trace context, in the reserved \texttt{traceparent} key. The specification
carves out an explicit exception to the reverse-DNS prefix rule for exactly this,
which tells you how important the working group considered it.

Without it, a slow request in a three-level hierarchy is eleven separate log
searches and a guess. With it, it is one waterfall. Chapter~\ref{ch:measuring}
builds the waterfall.

\section{Meridian's supervisor}

Three specialists, one coordinator, one traced request.

\begin{py}
supervisor = Host(name="meridian-supervisor",
                  consent=ConsentPolicy(auto_allow_read_only=True),
                  tool_token_budget=3000)

supervisor.connect("risk",       StreamableHttpClient(RISK_AGENT_URL))
supervisor.connect("compliance", StreamableHttpClient(COMPLIANCE_AGENT_URL))
supervisor.connect("fraud",      StreamableHttpClient(FRAUD_AGENT_URL))
supervisor.discover_all()
supervisor.refresh_catalogue()

result = AgentLoop(supervisor, model,
                   max_steps=6,
                   token_budget=40_000,
                   cost_budget_usd=0.25,
                   parallel_fanout=True).run(
    "Full review of ACC-1042 for the credit committee.")
\end{py}

\texttt{parallel\_fanout=True} is doing heavy lifting here. The three specialists
are independent, so the supervisor should call them together. From
Chapter~\ref{ch:primer}, that is a 2.8$\times$ speedup on the tool-call band, and
in a hierarchy each level's fan-out compounds.

\subsection{What the trace looks like}

\begin{plain}
supervisor.analyse_portfolio                        [ 4,180 ms]
 |- model turn (plan)                               [   440 ms]
 |- parallel fan-out                                [ 1,820 ms]
 |   |- risk-agent.analyse_account                  [ 1,820 ms]
 |   |   |- model turn (plan)                       [   430 ms]
 |   |   |- risk.assess_account_risk                [    12 ms]
 |   |   |- marketdata.price_facility               [     9 ms]
 |   |   `- model turn (synthesise)                 [   410 ms]
 |   |- fraud-agent.screen                          [   890 ms]
 |   `- compliance-agent.review                     [ 1,240 ms]
 `- model turn (synthesise)                         [   470 ms]
\end{plain}

Read the shape rather than the numbers. The fan-out costs the \emph{maximum} of
its children (1,820) rather than the sum (3,950), which is what parallelism
bought. And the leaf tool calls are 12\,ms and 9\,ms against model turns of 430
and 410, which is Chapter~\ref{ch:primer}'s point restated at a different scale:
you are paying for thinking, not for calling.

\keyidea{Every level of delegation adds at least two model turns: one to decide
what to delegate and one to read what came back. Delegate because the
decomposition genuinely helps, not because the architecture diagram looks
tidier.}

\section{When not to build a multi-agent system}

The section this chapter exists for.

Multi-agent architectures are intellectually appealing and expensive. Each
delegation adds model turns, adds a failure mode, adds a trace to correlate, and
adds a budget to enforce. Sometimes that is worth it. Often it is not.

\begin{table}[htbp]
\centering
\small
\begin{tabularx}{\textwidth}{@{}XX@{}}
\toprule
\thead{Delegate when} & \thead{Do not when} \\
\midrule
The sub-task needs a genuinely different tool catalogue
  & You could have added three tools to one server \\
The sub-task needs different authorization, and the boundary is real
  & The boundary is organisational rather than technical \\
Sub-tasks are independent and parallel fan-out pays
  & They are sequential, so you pay depth for no concurrency \\
Context isolation is the point: the child sees data the parent must not
  & You are isolating context you could have summarised in one turn \\
Different teams own different agents and deploy independently
  & One team owns all of it \\
\bottomrule
\end{tabularx}
\caption{The right-hand column describes most systems that have a multi-agent
architecture.}
\label{tab:when-delegate}
\end{table}

The honest test: \emph{would this be a function call if the components were in
one process?} If yes, and you have made it an agent, you have added several
hundred milliseconds and a distributed systems problem in exchange for an
organisational boundary you could have drawn with a module.

\section{Security across agent boundaries}

Every boundary is a trust boundary, and delegation makes it easy to forget that.

\begin{dangerbox}{Confused deputy, delegated}
A supervisor holds broad authorization. A specialist agent asks it to call a
tool. If the supervisor executes on behalf of the specialist using the
supervisor's own credentials, the specialist has just escalated to the
supervisor's privileges.

Which is the confused deputy from Chapter~\ref{ch:authorization}, one layer up
and much easier to build by accident, because the specialist is ``one of ours''.
\end{dangerbox}

Three rules:

\textbf{Delegate authorization, not just work.} The child should act with its
own credentials, scoped to what it needs. Token exchange exists for this.

\textbf{Never pass through a token.} It was forbidden in
Chapter~\ref{ch:authorization} and it is still forbidden here.

\textbf{Audit the delegation, not just the leaves.} An audit log showing
\texttt{assess\_account\_risk} was called tells you nothing about who decided to
call it. Record the chain.

And the one that people find surprising: an agent that consumes untrusted content
and then calls another agent has extended the injection surface. If a compliance
agent reads poisoned guidance text and then delegates to a risk agent, the
injection now reaches a system with different privileges.
Chapter~\ref{ch:threats} follows that thread to its conclusion.

\section{What to remember}

\begin{itemize}[leftmargin=1.6em, itemsep=3pt]
\item An agent is a host on one side and a server on the other. The protocol
needs no extension for this.
\item Supervisor trees for isolation, peer meshes for resilience, hierarchies for
context management. Pick for the property you need.
\item Four counters travel in \texttt{\_meta}: depth, token budget, cost budget,
trace id.
\item Check depth before doing work. Default to a small budget, never an
unlimited one.
\item Budgets must be remaining-and-decrementing, not per-node allowances, or
fan-out multiplies them.
\item Circuit breakers with a half-open probe, because a timeout burns the
caller's step budget and tells it nothing.
\item Propagate \texttt{traceparent}, or a slow request becomes eleven log
searches.
\item Every delegation adds at least two model turns. Delegate for real reasons.
\item Delegate authorization along with work, never a token.
\end{itemize}

That closes Part~III. You can build all of it now. Part~IV is about finding out
whether it is any good.
